\scanpage{10}
Or, si on part de $\varphi$ i.e de
\[
 f:E\longrightarrow E' \tag*{(36)}
\]
on peut lui associer
\[
 \bigwedge^2 f:\bigwedge^2 E\longrightarrow\bigwedge^2 E' \tag*{(37)}
\]
qui en vertu des isom. (6) définit une application linéaire
\[
 f'=\alpha'\circ\bigwedge^2 f\circ\alpha^{-1}:E'\longrightarrow E \tag*{(38)}
\]
\deleted{\unsure{Du point|De sorte}} Par définition, on aura
\[
 f'(\alpha(\deleted{\varphi}\unsure{x'|x}\wedge\unsure{y'|y}))=\alpha'(f(x)\wedge f(y))\qquad x,y\in\Gamma(S,E) \tag*{(39)}
\]
i.e.
\[
 \langle f'(\alpha(x\wedge y)),x'\rangle\omega'=f(x)\wedge f(y)\wedge x'\qquad(x'\in\Gamma(S,E')\unsure{)} \tag*{(40)}
\]
En fait, en termes d'une base \inserted{(8) et de sa duale}, les coefficients de la matrice \inserted{exprimant} $f'$ sont les 9 mineurs de la matrice de $f$ -- ce qui montre en particulier que si la forme \inserted{bil.} $\varphi$ associée à $f$ est symétrique, il en est de $=$ de la forme $\varphi'$ associée à $f'$. Dans ce cas, la structure d'algèbre de $E$ associée \deleted{$[\unsure{x,y|\, ,\,}]_{\varphi'}$} (qui n'est autre que
\[
 [x,y]_{\varphi'}=\alpha'\bigwedge^2 f(x\wedge y)=\alpha'(f(x)\wedge f(y)) \tag*{(40 bis)}
\]
\unsure{explicité|explicit.} par (40), qui s'écrit
\[
 \langle[x,y]_{\varphi'},x'\rangle\omega'=f(x)\wedge f(y)\wedge x'\qquad ) \tag*{(40 \deleted{bis}\inserted{ter})}
\]
est \inserted{encore} une structure d'\uline{algèbre de Lie}.

% archive sheet number: 9
